\documentclass[12pt, a4paper]{article}

\usepackage[utf8]{inputenc}
\usepackage{graphicx}
\usepackage{hyperref}
\usepackage{float}
\usepackage{geometry}
\geometry{a4paper, margin=1in}
\usepackage{titlesec}
\usepackage{enumitem}

\title{\textbf{DOKUMEN ANALISIS DAN PERANCANGAN SISTEM\\ Sistem Pengelolaan Nilai Akademik}}
\author{Dokumentasi Proyek}
\date{\today}

\begin{document}

\maketitle
\tableofcontents
\newpage

\section{Pendahuluan}
Dokumen ini berisi analisis dan perancangan untuk \textbf{Sistem Pengelolaan Nilai}. Sistem ini dirancang untuk memfasilitasi pengelolaan nilai siswa secara digital, transparan, dan efisien bagi sekolah, serta dibangun dengan menggunakan pendekatan pemrograman terstruktur dan \textit{Object-Oriented Programming (OOP)} pada *framework* Laravel.

\section{Tujuan Sistem}
Tujuan utama dari sistem pengelolaan nilai ini adalah:
\begin{itemize}
    \item Mendigitalisasi proses input dan rekapitulasi nilai siswa dari sistem manual menjadi terkomputerisasi.
    \item Memudahkan guru dalam memasukkan, mengedit, dan mengelola komponen nilai (Tugas, UTS, UAS).
    \item Memberikan transparansi kepada siswa untuk melihat hasil evaluasi akademik serta status kelulusan (Lulus/Tidak Lulus).
    \item Menghindari kehilangan data nilai melalui penyimpanan \textit{database} relasional yang aman.
\end{itemize}

\section{Analisis Kebutuhan Pengguna}
Berdasarkan studi kasus, terdapat tiga aktor utama dalam sistem:
\begin{enumerate}
    \item \textbf{Admin:} Membutuhkan antarmuka untuk mengelola pengguna (User) dan data referensi sistem.
    \item \textbf{Guru:} Membutuhkan fitur untuk mencatatkan nilai tugas, UTS, dan UAS terhadap siswa berdasarkan mata pelajaran yang diampu.
    \item \textbf{Siswa:} Membutuhkan halaman informasi atau rapor yang menampilkan rekapitulasi nilai mereka beserta status akhir.
\end{enumerate}

\section{Fungsi Utama Sistem}
Fungsi esensial yang harus ada pada sistem ini meliputi:
\begin{itemize}
    \item \textbf{Manajemen Autentikasi:} Mekanisme *Login/Logout* dengan perutean berbasis peran (Admin, Guru, Siswa).
    \item \textbf{Manajemen Data Master:} Sinkronisasi tabel profil guru dan profil siswa yang merujuk pada kredensial akun pengguna (*User Account*).
    \item \textbf{Manajemen Nilai:} Kalkulasi otomatis Nilai Akhir (Tugas, UTS, UAS) dan penentuan status kelulusan.
    \item \textbf{Pelaporan (Reporting):} Rekapitulasi nilai dan laporan PDF (Cetak Rapor).
\end{itemize}

\section{Spesifikasi Fungsional dan Non-Fungsional}
\subsection{Spesifikasi Fungsional}
\begin{itemize}
    \item Sistem dapat mengautentikasi pengguna dan menyimpan sesi (*session*) dengan aman.
    \item Sistem dapat memproses input tiga jenis nilai: \texttt{nilai\_tugas}, \texttt{nilai\_uts}, dan \texttt{nilai\_uas}.
    \item Sistem dapat menghitung \texttt{nilai\_akhir} berdasarkan bobot yang ditentukan.
    \item Sistem dapat mengekspor rekap nilai dalam bentuk PDF.
\end{itemize}

\subsection{Spesifikasi Non-Fungsional}
\begin{itemize}
    \item \textbf{Kinerja:} Kueri *database* dieksekusi secara optimal untuk mempercepat *loading* halaman.
    \item \textbf{Keamanan:} Sandi pengguna (Password) wajib di-\textit{hash} menggunakan fungsi \textit{Bcrypt}.
    \item \textbf{Antarmuka:} Bersifat responsif (*mobile-friendly*) dengan balutan *framework* CSS modern.
\end{itemize}

\section{Alur Kerja Sistem (Workflow)}
\begin{enumerate}
    \item Pengguna (\textit{Admin}, \textit{Guru}, atau \textit{Siswa}) melakukan masuk log (\textit{Login}).
    \item Guru memilih menu Input Nilai, memilih Siswa, lalu memasukkan Nilai Tugas, UTS, dan UAS untuk Mata Pelajaran tertentu.
    \item Sistem secara otomatis memanggil fungsi kalkulasi untuk mengisi kolom \textit{Nilai Akhir} dan \textit{Status}.
    \item Siswa login, mengakses dasbor nilai, lalu sistem menyajikan tabel capaian studinya.
\end{enumerate}

\section{Unified Modeling Language (UML) Sistem}
Diagram Use Case di bawah merepresentasikan secara visual fungsionalitas dan interaksi aktor.

\begin{figure}[H]
    \centering
    \includegraphics[width=0.8\textwidth]{diagram.png}
    \caption{Diagram Use Case Sistem Pengelolaan Nilai}
    \label{fig:usecase}
\end{figure}

\section{Rancangan Antarmuka (UI)}
Rancangan UI berfokus pada kemudahan pengguna (\textit{User Friendly}):
\begin{itemize}
    \item \textbf{Halaman Login:} Menangani autentikasi dengan \textit{form} input standar.
    \item \textbf{Dashboard (Admin/Guru/Siswa):} Berisi panel statistik, ringkasan, atau selamat datang (*Welcome Screen*).
    \item \textbf{Form Input Nilai (Guru):} *View Blade* yang berisi kotak masukan (\textit{input textbox}) untuk mencatat poin tugas, UTS, dan UAS.
    \item \textbf{Halaman Laporan (Siswa):} *Grid* tabel nilai yang dapat dicetak menjadi lampiran fisik (PDF).
\end{itemize}

\section{Rancangan Database (ERD)}
Struktur relasional \textit{database} dirancang sesuai dengan migrasi Laravel terbaru proyek ini.

\begin{figure}[H]
    \centering
    \includegraphics[width=0.7\textwidth]{erd.png}
    \caption{Entity Relationship Diagram (ERD)}
    \label{fig:erd}
\end{figure}

Keterangan struktur tabel aktual:
\begin{itemize}
    \item \textbf{users:} Mengelola hak akses dan kata sandi dengan atribut \texttt{role} berjenis \texttt{enum('admin', 'guru', 'siswa')}.
    \item \textbf{siswas:} Menampung identitas unik siswa (\texttt{nis, nama, kelas}) dan berelasi One-to-One dengan tabel \texttt{users}.
    \item \textbf{gurus:} Menampung identitas pengajar (\texttt{kode\_guru, nama, mata\_pelajaran}) dan berelasi One-to-One dengan tabel \texttt{users}.
    \item \textbf{nilais:} Tabel transaksional. Berisi \textit{Foreign Key} (\texttt{siswa\_id, guru\_id}), kolom subjek (\texttt{mata\_pelajaran}), skor per komponen (\texttt{nilai\_tugas, nilai\_uts, nilai\_uas, nilai\_akhir}), dan predikat capaian (\texttt{status}).
\end{itemize}

\section{Batasan Sistem}
Untuk menjaga fokus aplikasi, batasan ruang lingkupnya adalah:
\begin{itemize}
    \item Sistem tidak mencakup pendaftaran Penerimaan Peserta Didik Baru (PPDB).
    \item Tidak memiliki integrasi pembayaran (*Payment Gateway*).
\end{itemize}

\section{Rancangan Pemrograman \& Status Database}

\subsection{Koneksi Database Berhasil}
\textbf{Status Koneksi Database: BERHASIL.} \\
Aplikasi telah terhubung dengan \textit{database} MySQL. Migrasi skema (*schema migrations*) untuk tabel \texttt{users}, \texttt{siswas}, \texttt{gurus}, dan \texttt{nilais} telah dieksekusi sempurna pada mesin lokal tanpa galat koneksi basis data.

\subsection{Rancangan Fungsi / Prosedur (Implementasi Pemrograman Terstruktur)}
Sistem ini menggunakan minimal \textbf{3 fungsi terstruktur / blok operasional} independen:
\begin{enumerate}
    \item \textbf{Fungsi Kalkulasi \texttt{hitungNilaiAkhir(\$tugas, \$uts, \$uas)}} \\
    \textit{Logika:} Menyatukan rasio bobot persentase dari parameter yang dimasukkan, misalnya \texttt{(30\% * Tugas) + (30\% * UTS) + (40\% * UAS)} untuk menghasilkan desimal (*float*) nilai akhir murni.

    \item \textbf{Fungsi Penentuan Status \texttt{tentukanStatus(\$nilaiAkhir)}} \\
    \textit{Logika:} Prosedur pengkondisian IF-ELSE. Jika variabel \texttt{\$nilaiAkhir >= 75}, prosedural ini akan membalikkan parameter konstan \textit{String} \texttt{"Lulus"}. Apabila di bawah 75, maka \texttt{"Tidak Lulus"}.

    \item \textbf{Fungsi Agregasi Rekap \texttt{hitungRataRataKelas(\$koleksiNilai)}} \\
    \textit{Logika:} Melakukan iterasi *looping* (misalnya \texttt{foreach} atau \texttt{array\_reduce}) pada kumpulan (koleksi) model nilai, dan menjumlahkannya lalu membaginya dengan total pencacah (*counter*).
\end{enumerate}

\subsection{Rancangan Class / Method (Implementasi OOP)}
Implementasi *Object-Oriented Programming* ditunjukkan kuat pada lapisan *Model* dan *Controller*:

\begin{enumerate}
    \item \textbf{Class \texttt{Nilai} (Model)}
    \begin{itemize}
        \item \textbf{\texttt{public function hitungNilaiAkhir()}}: \\
        \textit{Method OOP internal} yang diaktifkan (*triggered*) melalui \textit{booted lifecycle hook} (seperti \texttt{static::saving}) pada model Laravel. Mengambil properti di dalam objek (\texttt{\$this->nilai\_tugas}, dsb) tanpa harus menerima parameter eksternal.
        \item \textbf{\texttt{public function tentukanStatus()}}: \\
        \textit{Method} utilitas spesifik objek \texttt{Nilai}. Mengotomatisasi pembaruan nilai kolom properti \texttt{status} ke 'Lulus' atau 'Tidak Lulus' sebelum kueri SQL INSERT dijalankan.
    \end{itemize}

    \item \textbf{Class \texttt{GuruController} (Controller)}
    \begin{itemize}
        \item \textbf{\texttt{public function storeNilai(Request \$request)}}: \\
        \textit{Method} ini bertugas menangkap *payload* data *request*, memvalidasinya, kemudian memanggil metode penyimpanan *Object-Relational Mapping* (ORM) ke \textit{database}. Mengembalikan sinyal \textit{RedirectResponse} atau \textit{JsonResponse}.
        \item \textbf{\texttt{public function dashboard()}}: \\
        \textit{Method} ini merangkum data kontekstual untuk peran Guru, lalu merendernya dengan *layout engine* \textit{Blade View} Laravel.
    \end{itemize}
\end{enumerate}

\end{document}
